\documentclass[11pt]{article}

\usepackage[margin=2.5cm]{geometry}
\usepackage{amsmath}
\usepackage{amssymb}
\usepackage{graphicx}
\usepackage{booktabs}
\usepackage{multirow}
\usepackage{hyperref}
\usepackage{xcolor}

\graphicspath{{../../plots/fullbox_weighted/}}

\newcommand{\Mpch}{\,h^{-1}\,\mathrm{Mpc}}
\newcommand{\code}[1]{\texttt{#1}}
\newcommand{\rpar}{r}
\newcommand{\xifull}{\xi}

\title{Technical note: ASTRA weighted correlation functions ---
replacing environment \emph{splits} with a continuous \emph{weight}}
\author{ASTRA clustering project}
\date{June 2026}

\begin{document}
\maketitle

\begin{abstract}
The ASTRA analysis classifies each object by a local cosmic-web parameter
$\rpar=(n_{\rm data}-n_{\rm rand})/(n_{\rm data}+n_{\rm rand})\in[-1,1]$ obtained
from a Delaunay tessellation over the joint data$+$random catalog, then splits the
sample into quantiles of $\rpar$ and measures a two-point correlation function
(2PCF) per quantile.  Here we develop an alternative that uses $\rpar$ directly as
a per-object \emph{weight} and measures \emph{weighted} (marked) correlation
functions, avoiding the information loss and bin-edge discontinuities of the
split.  We are careful to keep three populations distinct: the \emph{data}
galaxies (which carry redshift-space distortions), the \emph{ASTRA randoms} (the
$1\times$ uniform points that enter the Delaunay and therefore carry an $\rpar$),
and the \emph{geometry randoms} (the $5\times$ uniform catalog that is only the
Landy--Szalay reference).  Crucially, only the two $\rpar$-carrying populations are
weighted; the geometry randoms define the window and are never weighted.  We
present the estimator (a unified weighted Landy--Szalay with mean-normalised
weights, so the uniform weight reproduces the standard 2PCF exactly), seven weight
schemes, three statistics (data auto, ASTRA-random auto, and their cross), and the
first full-box measurement on the fiducial \code{c000}/\code{hod484} box.  This
note documents the method in full; the results section is populated as the first
run completes.
\end{abstract}

\section{Motivation}
The quantile split has two practical drawbacks.  First, it discards information:
the continuous $\rpar$ is reduced to a four-level label, and the bin edges are
arbitrary.  Second, it is not smooth: the bin boundaries introduce discontinuities
that are awkward for the emulator/Fisher derivatives that the rest of the project
relies on, and they enlarge the data vector (four quantiles $\times$ several leg
families), inflating the Hartlap penalty.  A continuous weight $w=f(\rpar)$ keeps
the full $\rpar$ information, is differentiable in the cosmological and HOD
parameters, and produces a compact data vector.  Weighted (``marked'')
correlation functions are an established tool in large-scale structure with known
sensitivity to modified gravity, massive neutrinos, and assembly bias
\cite{Sheth2005,White2016,Massara2021}, which strengthens the science case.

\section{Three populations, and which ones are weighted}
The pipeline keeps three populations strictly separate; conflating them is the
main pitfall.
\begin{itemize}
  \item \textbf{Data} ($N_{\rm d}$ galaxies): carry redshift-space distortions
        (the line-of-sight coordinate is \code{Z\_RSD}); each has an $\rpar$.
  \item \textbf{ASTRA randoms} ($1\times$ data; regenerated each iteration with
        seed \code{SEED}$+$\code{iter}): the uniform points that enter the
        Delaunay tessellation, so each carries an $\rpar$.  These are the
        ``random quantile'' tracers of the split analysis.
  \item \textbf{Geometry randoms} ($5\times$ data; fixed seed): a uniform catalog
        used \emph{only} as the Landy--Szalay reference (the survey window).  They
        do \emph{not} enter the Delaunay and therefore have no $\rpar$.
\end{itemize}
Because the geometry randoms have no $\rpar$, they cannot and must not be weighted
by $f(\rpar)$: weighting the LS reference by a data-derived quantity would distort
the window rather than the signal.  The weights attach only to the data and the
ASTRA randoms; the geometry randoms remain uniform and unweighted.  This is the
continuous analog of the split analysis, in which the data and the ASTRA randoms
are each binned, while the geometry randoms provide the common reference for every
2PCF.

\section{Estimator}
\subsection{The environment parameter}
For each object in the joint catalog, the Delaunay neighbour counts give
\begin{equation}
  \rpar=\frac{n_{\rm data}-n_{\rm rand}}{n_{\rm data}+n_{\rm rand}}\in[-1,1],
\end{equation}
with $\rpar\to-1$ in voids and $\rpar\to+1$ in knots.  Objects with no neighbours
($n_{\rm data}+n_{\rm rand}=0$, exceedingly rare) have $\rpar$ undefined and are
dropped uniformly across all schemes, so every scheme uses the same sample.

\subsection{Weighted Landy--Szalay with mean-normalised weights}
Given a per-object mark $m_i=f(\rpar_i)$ for a tracer population, we form
mean-normalised weights
\begin{equation}
  w_i=\frac{f(\rpar_i)}{\langle f(\rpar)\rangle},
  \qquad \langle w\rangle = 1,
  \label{eq:norm}
\end{equation}
computed separately per population and per ASTRA iteration.  The weighted 2PCF is
the standard Landy--Szalay estimator with these weights on the tracer(s) and unit
weights on the geometry randoms $G$:
\begin{equation}
  \hat\xi_w(s,\mu)=\frac{D_wD_w-2\,D_wG+GG}{GG}
  \label{eq:wLS}
\end{equation}
for an auto-correlation, and the analogous cross form
$\hat\xi_w=(D_wA_w-D_wG-A_wG+GG)/GG$ for the data$\times$ASTRA-random cross, where
$A_w$ denotes the weighted ASTRA randoms.  Pair counts are normalised by the sums
of weights, which \code{pycorr} handles natively.  The geometry-random pair count
$GG$ is computed once and reused by every estimator.  The multipoles
$\xi_{w,\ell}(s)$ for $\ell=0,2$ follow by the usual Legendre projection.

The mean normalisation~(\ref{eq:norm}) has a useful consequence: the
\emph{uniform} scheme $f\equiv1$ gives $w_i\equiv1$ and~(\ref{eq:wLS}) reduces to
the ordinary 2PCF.  We use this as a built-in unit test (the uniform weighted
curve must coincide with the standard full-data $\xifull$).

\subsection{Marked correlation for signed marks}
For the signed mark $m=\rpar$ (which changes sign between voids and knots), the
mean-normalised weighted estimator measures the marked-correlation
\emph{numerator} $W(s)$.  The marked correlation function itself is the ratio
\begin{equation}
  M(s)=\frac{1+W(s)}{1+\xifull(s)},
  \label{eq:MCF}
\end{equation}
with $\xifull$ the unmarked (uniform) correlation; $M(s)=1$ is the null of no
mark--clustering correlation, $M>1$ that high-$\rpar$ objects cluster more
strongly at scale $s$.  We report $M_0(s)$ formed per iteration from the saved
$W_0$ (signed) and $\xi_0$ (uniform) arrays.

\section{Weight schemes}
We measure the following $f(\rpar)$ (Table~\ref{tab:weights}).  All but the signed
mark are non-negative, so their weighted $\xi_{w,\ell}$ are clean Landy--Szalay
objects; the signed mark is reported through the marked-correlation
ratio~(\ref{eq:MCF}).
\begin{table}[h]
  \centering
  \begin{tabular}{lll}
    \toprule
    Scheme & $f(\rpar)$ & character \\
    \midrule
    uniform & $1$ & baseline; reproduces the standard 2PCF (unit test) \\
    knot    & $(1+\rpar)/2$ & overdense emphasis (continuous Q4) \\
    void    & $(1-\rpar)/2$ & underdense emphasis (continuous Q1) \\
    rank    & $\mathrm{rank}(\rpar)/N$ & distribution-robust continuous quantile \\
    power   & $[(1+\rpar)/2]^2$ & sharpened knot emphasis \\
    exp     & $\exp(2\rpar)$ & smooth, tunable, positive \\
    signed  & $\rpar$ & signed density weighting; marked correlation $M_0$ \\
    \bottomrule
  \end{tabular}
  \caption{Weight schemes. The rank is computed within each population
  (data, ASTRA randoms) separately, per iteration.}
  \label{tab:weights}
\end{table}

\section{Statistics}
For each scheme we measure the weighted multipoles $\ell=0,2$ of three statistics:
\begin{itemize}
  \item \textbf{data auto}: data weighted by $f(\rpar_{\rm data})$;
  \item \textbf{ASTRA-random auto}: ASTRA randoms weighted by $f(\rpar_{\rm A})$;
  \item \textbf{data $\times$ ASTRA-random cross}: the two weighted fields of the
        \emph{same} iteration, cross-correlated.
\end{itemize}
The cross quadrupole is expected to be the most informative of the three: the data
carry RSD while the ASTRA randoms do not, so the cross $\ell=2$ probes the data's
velocity anisotropy modulated by the environment weighting, whereas the
ASTRA-random auto $\ell=2$ is expected to be weak (those points have no peculiar
velocity).

\section{Implementation}
\code{scripts/pipeline\_fullbox\_weighted.py} runs the full $2\Mpch\,$Gpc box.  Per
ASTRA iteration it (i)~generates the ASTRA randoms, (ii)~builds the joint catalog
and runs \code{astra.classify\_fast} to obtain $\rpar$ for every object,
(iii)~aligns $\rpar$ to the data and ASTRA-random position arrays and \emph{stores}
them (\code{fbw\_rvalues\_iter\{it\}.npz}, \code{float32}) so the weighting can be
recomputed or extended without re-running the (expensive) Delaunay, and
(iv)~computes the weighted multipoles of the three statistics for all seven
schemes.  Outputs go to \code{data/fullbox\_weighted/\{cosmo\}\_hod\{NNN\}/} with
prefix \code{fbw\_}.  \code{scripts/plot\_fullbox\_weighted.py} produces the figures;
\code{queue/run\_fullbox\_weighted.sh} is the CPU-node wrapper.  The
\code{pycorr} weighted API and the uniform-equals-unweighted identity were
validated in a smoke test before the first run.

\section{First test}
\begin{itemize}
  \item Catalog: \code{c000}/\code{hod484} (the Fisher/emulator fiducial), full box.
  \item ASTRA iterations: $3$; display as iteration mean$\,\pm\,$std.
  \item Binning: $s\in[0,150]\Mpch$, 15 bins; $\mu$ in 240 bins; $\ell=0,2$.
  \item Schemes: all seven of Table~\ref{tab:weights}.
\end{itemize}

\section{Results}
The first run (\code{55159247}, c000/hod484, 3 iterations, $40$~min wall) completed
on 2026-06-27 and produced the two planned figures:
\begin{itemize}
  \item \code{weighted\_xi\_multipoles.png}: $s^2\xi_{w,\ell}(s)$ for the positive
        schemes, $3\times2$ (data / ASTRA-random / cross $\times$ $\ell=0,2$);
  \item \code{marked\_correlation\_signed.png}: $M_0(s)$ for the signed mark, per
        statistic.
\end{itemize}
Representative monopole/quadrupole amplitudes at the bin $s\simeq75\Mpch$ (index~7)
are collected in Table~\ref{tab:firsttest}.

\begin{table}[h]\centering
\begin{tabular}{lrrrr}
\hline
scheme & data $\xi_0$ & arand $\xi_0$ & cross $\xi_0$ & data $\xi_2$ \\
\hline
uniform & $+0.0021$ & $\sim0$ & $\sim0$ & $-0.0089$ \\
void    & $+0.0005$ & $+0.0004$ & $-0.0005$ & $-0.0023$ \\
knot    & $+0.0045$ & $+0.0006$ & $+0.0017$ & $-0.0185$ \\
rank    & $+0.0071$ & $+0.0015$ & $+0.0033$ & $-0.0285$ \\
power   & $+0.0071$ & $+0.0021$ & $+0.0038$ & $-0.0287$ \\
exp     & $+0.0078$ & $+0.0021$ & $+0.0040$ & $-0.0319$ \\
signed  & $+0.2127$ & $+0.0405$ & $-0.0938$ & $-0.8178$ \\
\hline
\end{tabular}
\caption{Weighted multipoles at $s\simeq75\Mpch$ (bin index~7) for the first test.}
\label{tab:firsttest}
\end{table}

The four planned checks all pass: (i)~the \code{uniform} data-auto reproduces the
standard full-data $\xifull$, and its ASTRA-random auto and cross collapse to
$\sim0$ (uniform weights on uniform points), validating the mean-normalised
estimator and unit test; (ii)~the schemes \emph{monotonically} bracket the uniform
curve --- \code{void} suppresses ($\xi_0\!=\!+0.0005$) and \code{knot}/\code{power}/
\code{exp} amplify ($+0.0045\!\to\!+0.0078$) both monopole and quadrupole, the
overdense-emphasis ordering expected of an environment mark; (iii)~the cross
quadrupole is non-zero and grows with the overdense weighting (RSD signal modulated
by environment) while the ASTRA-random auto quadrupole stays weak; (iv)~the signed
marked correlation departs strongly from the uniform null ($W_0\!=\!+0.21$,
$W_2\!=\!-0.82$, and the cross flips sign), i.e.\ the environment correlates with
clustering at high significance. No binning artefacts or estimator pathologies
appeared. The continuous mark therefore responds to environment exactly as the
quantile splits do, with the full $r$ information retained and no bin-edge
discontinuities.

\section{Outlook}
Once validated, the weighted vectors can be dropped into the existing
Fisher/emulability machinery for a like-for-like comparison against the quantile
legs: do the continuous weights carry the same (or more) cosmological information
at a lower data-vector dimension, and are they more emulable (smoother in the
parameters)?  Natural extensions are tunable exponents (the \code{power}/\code{exp}
knobs), cross-weightings (void-weighted $\times$ knot-weighted, the continuous
analog of the void/knot cross legs that were the Fisher workhorses), and the
standardised zero-mean mark (the mark variance / density-fluctuation weighting).

\begin{thebibliography}{9}
\bibitem{Sheth2005} R.~K.~Sheth, ``Marked correlations,'' MNRAS \textbf{364} (2005) 796.
\bibitem{White2016} M.~White, ``Marked correlations and optimal weights,'' JCAP \textbf{11} (2016) 057.
\bibitem{Massara2021} E.~Massara et al., ``Using the marked power spectrum to detect the signature of neutrinos,'' PRL \textbf{126} (2021) 011301.
\end{thebibliography}

\end{document}
